\section{Continuum Derivation Details}
\label{app:continuum-derivation-details}

This appendix supplements the continuum mechanics details that support the main-body
assumptions without interrupting the reduced-order forward-model
narrative within this report. It is appendix background for notation and
model hierarchy, not additional evidence for the first controlled 1D study.
All identities below are classical identities under the regularity needed for
the displayed derivatives, traces, and changes of variables. The standing
regularity convention, boundary notation, and differential-operator convention
are those of
Convention~\ref{conv:foundation-standing-regularity},
Definition~\ref{def:foundation-domain-boundary-notation}, and
Definition~\ref{def:foundation-differential-operator-conventions}. In
particular, \(\Divbf\) is the row-divergence of a tensor field, and boundary
forms assume the stated trace and normal hypotheses; no weak formulation is
being asserted in this appendix.

\subsection{Material Derivative and Trajectory Details}
\label{subsec:material-derivative-trajectory-details}

For a scalar field $\phi\in C^1(\mathcal{Q}_T)$ and a sufficiently regular flow
map \(\La_t\), define the Lagrangian pullback
$\phihat(t,\vxi):=\phi(t,\La_t(\vxi))$. Differentiating with respect to time
along a trajectory and using
$\pd_t \La_t(\vxi)=\vu(t,\La_t(\vxi))$ gives
\begin{flalign*}
    \quad \dt\phi(t,\La_t(\vxi)) = \pd_t \phi(t,\vx)+\nabla\phi(t,\vx)\cdot\vu(t,\vx), \qquad \vx=\La_t(\vxi). &&
\end{flalign*}
The identity
\begin{flalign*}
    \quad \Div(\phi\vu)=\nabla\phi\cdot\vu+\phi\Div\vu &&
\end{flalign*}
also gives
\begin{flalign*}
    \quad D_t\phi=\pd_t \phi+\Div(\phi\vu)-\phi\Div\vu. &&
\end{flalign*}

\subsection{Jacobian Evolution and Reynolds Transport}
\label{subsec:jacobian-evolution-reynolds-transport}

Let $\vFhat_t=\nabla_{\vxi}\La_t$ and $\Jhatdet_t=\det\vFhat_t$. Here \(\La_t\)
is assumed to be an orientation-preserving \(C^1\) change of variables on the
material volume being considered, with the additional time regularity needed to
differentiate \(\vFhat_t\) and \(\Jhatdet_t\). The standard Jacobi identity
gives
\begin{flalign*}
    \quad
    \pd_t \Jhatdet_t
    =
    \Jhatdet_t\,\mathrm{tr}\!\left(\pd_t \vFhat_t\,\vFhat_t^{-1}\right)
    =
    \Jhatdet_t\,\Div\vu(t,\La_t(\vxi)),
    &&
\end{flalign*}
which is the Lagrangian form of $D_t\Jdet_t=\Jdet_t\Div\vu$
\cite[Part~I, Ch.~1, Sec.~1.3]{hemodynamics}.

For $V(t)=\La_t(V(0))$, change variables in
$I(t)=\int_{V(t)}\phi(t,\vx)\dx$:
\begin{flalign*}
    \quad I(t) = \int_{V(0)}\phi(t,\La_t(\vxi))\Jhatdet_t(\vxi)\dxi. &&
\end{flalign*}
Differentiating under the integral sign and substituting the Jacobian-evolution
identity gives
\begin{flalign*}
    \quad
    I'(t)
    =
    \int_{V(t)}
    \left(D_t\phi+\phi\Div\vu\right)\dx
    =
    \int_{V(t)}
    \left(\pd_t \phi+\Div(\phi\vu)\right)\dx. &&
\end{flalign*}
If $V(t)$ has Lipschitz boundary and the trace of \(\phi\vu\) is integrable on
\(\partial V(t)\), the divergence theorem gives the boundary form
\begin{flalign*}
    \quad
    I'(t)=\int_{V(t)} \pd_t \phi\,\dx
    +\int_{\partial V(t)}\phi\,\vu\cdot\nhat\,dS. &&
\end{flalign*}

\subsection{Linear Momentum Balance Details}
\label{subsec:linear-momentum-balance-details}

Starting from the material-volume balance
\begin{flalign*}
    \quad
    \dt\int_{V(t)}\rho\vu\,dV
    =
    \int_{\partial V(t)}\vT\nhat\,dS
    +\int_{V(t)}\rho\vfb\,dV, &&
\end{flalign*}
Reynolds transport gives
\begin{flalign*}
    \quad
    \dt\int_{V(t)}\rho\vu\,dV
    =
    \int_{V(t)}
    \left(D_t(\rho\vu)+\rho\vu\,\Div\vu\right)dV. &&
\end{flalign*}
Using mass conservation,
$D_t\rho+\rho\Div\vu=0$, this reduces to
\begin{flalign*}
    \quad
    \int_{V(t)}\rho D_t\vu\,dV
    =
    \int_{V(t)}\rho\left(\pd_t \vu+(\vu\cdot\nabla)\vu\right)dV. &&
\end{flalign*}
The surface traction term satisfies
\begin{flalign*}
    \quad
    \int_{\partial V(t)}\vT\nhat\,dS
    =
    \int_{V(t)}\Divbf(\vT)\,dV. &&
\end{flalign*}
This is the tensor row-divergence convention from
Definition~\ref{def:foundation-differential-operator-conventions}, applied
under the regularity needed for the stress trace \(\vT\nhat\) and the volume
integral of \(\Divbf\vT\). Since $V(t)$ is arbitrary within the localization
class, the differential form is
\begin{flalign*}
    \quad
    \rho\left(\pd_t \vu+(\vu\cdot\nabla)\vu\right)
    =
    \Divbf(\vT)+\rho\vfb. &&
\end{flalign*}
The conservative Eulerian form
\begin{flalign*}
    \quad
    \pd_t (\rho\vu)+\Divbf(\rho\vu\otimes\vu)
    =
    \Divbf(\vT)+\rho\vfb &&
\end{flalign*}
is equivalent whenever the continuity equation holds, because
\begin{flalign*}
    \quad
    \pd_t (\rho\vu)+\Divbf(\rho\vu\otimes\vu)
    =
    \rho\left(\pd_t \vu+(\vu\cdot\nabla)\vu\right)
    +\vu\left(\pd_t \rho+\Div(\rho\vu)\right). &&
\end{flalign*}
